\chapter{Conclusão}

Este trabalho propôs, implementou e avaliou uma matheurística baseada em Large Neighborhood Search (LNS) com reotimização por sub-MIP, em um esquema do tipo \textit{fix-and-optimize}, integrada como heurística primal ao solver SCIP 7.0, para o Problema da Distribuição de Disciplinas (PDD). O problema foi modelado como um programa linear inteiro binário, cuja função objetivo maximiza a qualidade da atribuição de disciplinas a professores, respeitando restrições de carga horária mínima anual e máxima semestral.

A calibração dos parâmetros internos da LNS foi conduzida via abordagem \emph{one-factor-at-a-time}, avaliando três parâmetros: taxa de destruição, sentido de ordenação dos candidatos e tempo máximo do sub-MIP. Em todas as etapas, a métrica de nós restantes na árvore de Branch-and-Bound foi o critério decisivo, uma vez que tempo de execução e gap de otimalidade não apresentaram diferenças significativas entre os cenários avaliados. Para as instâncias testadas, a combinação de parâmetros selecionada --- 75\% de destruição, ordenação crescente por \texttt{avgPreferenceWeight} e 200 segundos de sub-MIP --- esteve associada a incumbentes de melhor qualidade e a uma convergência mais rápida do Branch-and-Bound.

No comparativo final, foram avaliadas quatro abordagens: B\&B puro, LNS, GRASP e a variante híbrida GRASP|LNS. O conjunto experimental reuniu todas as 42 instâncias sintéticas e uma instância real. Os resultados não indicaram uma abordagem dominante em todas as métricas. O B\&B puro obteve o menor tempo total (5.232,81s) e o menor gap médio (0,3473\%). A variante GRASP|LNS apresentou o menor total de nós restantes (187.673), uma redução de 54,9\% em relação ao B\&B puro. B\&B puro e LNS empataram no número de instâncias resolvidas até a otimalidade, com 32 de 43. O GRASP isolado apresentou o maior tempo total, o maior gap médio e o menor número de instâncias ótimas.

Na instância real, composta por 224 disciplinas e 61 professores, LNS e B\&B puro provaram a otimalidade e encerraram a árvore de busca. O LNS foi 2,93s mais rápido, enquanto GRASP e a variante híbrida GRASP|LNS atingiram o limite temporal com gaps positivos. Esse resultado demonstra desempenho competitivo do LNS no caso real avaliado, mas uma única instância não permite generalizar essa vantagem para outras instituições ou configurações do PDD.

Os resultados obtidos mostram que a integração de heurísticas ao Branch-and-Bound produz efeitos distintos conforme a métrica observada. A LNS não superou globalmente o método exato, mas permaneceu competitiva na quantidade de soluções ótimas e na instância real. A variante híbrida mostrou capacidade de reduzir a árvore remanescente, embora essa redução não tenha resultado no menor tempo ou gap médio. O trabalho cumpriu os objetivos propostos: modelagem do problema como PLI, implementação de uma heurística LNS do tipo \textit{fix-and-optimize} no SCIP, análise do impacto dos parâmetros internos e comparação abrangente com abordagens alternativas.

\section{Limitações}

As principais limitações deste trabalho são:

\begin{enumerate}
  \item A calibração \emph{one-factor-at-a-time} ajusta os parâmetros de forma sequencial, sem explorar possíveis interações entre eles. Uma análise fatorial completa ou métodos de \emph{design of experiments} poderiam revelar combinações mais eficientes.

  \item A avaliação incluiu apenas uma instância real. Embora ela complemente as 42 instâncias sintéticas e permita observar o comportamento das abordagens em um caso aplicado, não representa a diversidade de regras, preferências e dimensões encontrada entre instituições de ensino.

  \item A variante do GRASP utilizada na comparação (\texttt{max\_iter=1}, sem busca local) representa uma versão simplificada da heurística. Uma parametrização mais elaborada poderia alterar o desempenho relativo.

  \item O limite de tempo global fixo de 400 segundos por instância não permite avaliar o comportamento das abordagens sob restrições temporais mais rígidas ou mais flexíveis.
\end{enumerate}

\section{Trabalhos Futuros}

Como direções para trabalhos futuros, destacam-se:

\begin{enumerate}
  \item Investigar estratégias adaptativas para os parâmetros do LNS, ajustando dinamicamente a taxa de destruição e o tempo do sub-MIP com base no progresso da busca.

  \item Avaliar a matheurística em múltiplas instâncias reais, provenientes de diferentes instituições e períodos acadêmicos, incluindo problemas de maior porte e restrições adicionais como preferências de horário e indisponibilidades.

  \item Explorar critérios alternativos de seleção de candidatos à destruição, como métricas baseadas na contribuição marginal de cada atribuição à função objetivo.

  \item Investigar mecanismos de equidade e de resistência à declaração estratégica de preferências. Como a maximização da qualidade total pode concentrar ganhos em parte dos professores, um docente poderia declarar preferência por poucas disciplinas para tentar aumentar a prioridade de suas escolhas. Formulações multiobjetivo ou do tipo max-min, combinadas com informações como qualificação, histórico de atribuições ou limites mínimos de opções aceitáveis, podem reduzir essa vantagem e equilibrar a satisfação entre os professores.

  \item Comparar o desempenho com solvers comerciais (Gurobi, CPLEX) e com outras metaheurísticas, como Busca Tabu e Algoritmos Genéticos.

\end{enumerate}
